\section{Espaços de Sobolev}

\subsection*{Espaços de Sobolev}
\begin{frame}{Espaços de Sobolev}
\begin{defin}
 Sejam $\Omega \subset \R^\dm$, {\color{blue}$k\geq 0$} um inteiro e {\color{red}$1\leq p\leq +\infty$}. O \dt{Espaço de Sobolev} $W^{{\color{blue}k},{\color{red}p}}(\Omega)$  é o espaço vetorial normado definido por 
 \[W^{{\color{blue}k},{\color{red}p}}(\Omega):=\{u\in L^{\color{red}p}(\Omega);\ D^{\color{blue}\alpha} u\in L^{\color{red}p}(\Omega),\ \forall\, {\color{blue}|\alpha|\leq k}\},\]
 com a norma dada por 
 \[
 \n{u}_{{\color{blue}k},{\color{red}p},\Omega}=
\begin{cases}
\displaystyle\left(\sum_{{\color{blue}|\alpha|\leq k}} \n{D^{\color{blue}\alpha} u}^{\color{red}p}_{L^{\color{red}p}(\Omega)}\right)^{1/{\color{red}p}}, & \text{ se } {\color{red}p\in [1,+\infty)}\\[2em]
\displaystyle\sum_{{\color{blue}|\alpha|\leq k}} \n{D^{\color{blue}\alpha} u}_{L^{\color{red}\infty}(\Omega)}, & \text{ se } {\color{red}p=\infty}
\end{cases}
\]

\end{defin}



\end{frame}

\begin{frame}
	\begin{block}{Notações}
		\begin{enumerate}
			
			\item Note que $W^{0,p}(\Omega)=L^p(\Omega)$ e  $\n{u}_{0,p,\Omega}=\n{u}_{L^p(\Omega)}$ que denotaremos simplesmente por $\n{u}_{p,\Omega}$.
			
			\item Quando não houver perigo de confusão omitiremos $\Omega$ na notação da norma.
			
			
		\item   Se $p=2$, escrevemos $H^k(\Omega)=W^{k,2}(\Omega)$. 
		\end{enumerate}
	\end{block}

	\begin{exer}
	Mostre que $W^{k,p}(\Omega)$ é um espaço vetorial normado e que  $\n{u}_{k,p,\Omega}$ é equivalente à norma
	\[\n{u}=\sum_{|\alpha|\leq k}\n{D^\alpha u}_{L^p(\Omega)}.\]
\end{exer}
\end{frame}

\begin{frame}
	\begin{exe}
\begin{enumerate}[a]
	\item Sejam $\Omega=(-1,1)$ e	$u(x)=x\chi_{(0,1)}$. Note que $u\in L^1(-1,1)$, $u'\in L^1(-1,1)$ no sentido clássico, mas $Du=\delta_0\not\in L^1(-1,1)$ (não possui derivada fraca), portanto $u\not\in W^{1,1}(-1,1)$.
	
	\item Sejam $\Omega=B(0,1)$, $\alpha>0$ e $u(x)=\|x\|^{-\alpha}$, $x\in \Omega$, $x\neq 0$. Mostre que $u\in W^{1,p}(\Omega)$ se, e somente se, $0<\alpha<\frac{N-p}{p}$.
\end{enumerate}

	\end{exe}

\subsection*{Exercício \theexer}
\end{frame}


\begin{frame}
	\begin{prop}
	Para todo $1\leq p\leq +\infty$, $W^{k,p}(\Omega)$ é um espaço de Banach. Se $1<p<+\infty$ ele é reflexivo e se $1\leq p<+\infty$ é separável.
\end{prop}

\begin{lema}
	Sejam $\Omega$ um aberto de $\R^\dm$ e $1\leq p\leq +\infty$. A aplicação 
	\[u\in W^{1,p}(\Omega)\longmapsto (u,\partial_1 u,\ldots,\partial_\dm u)\in L^p(\Omega)^{\dm+1}\]
	é uma \textcolor{red}{ isometria}, onde estamos munindo o espaço $L^p(\Omega)^{\dm+1}$ com a norma
	\[\n{v}=\left(\sum_{i=1}^{\dm+1}\n{v_i}_{p,\Omega}^p\right)^{1/p}\]
\end{lema}
\end{frame}

%%%%%%%%%%%%%%%%%%%%%%%%%%%%%%%%%%%%%%%%%%%%%%%%%%%%%%%%%%%%%%%%%%%%%%
\subsection*{O espaço \texorpdfstring{\(W_0^{1,p}(\Omega)\)}{W₀¹ 𞀾(Ω)}}

\begin{frame}{Espaços $W_0^{k,p}(\Omega)$}
Como vimos na Proposição \ref{densLp}, $\D(\Omega)$ é denso em $L^p(\Omega)$, mas isto não é verdade para $W^{k,p}(\Omega)$. Por essa razão, define-se o espaço 
\[W_0^{k,p}(\Omega):=\overline{\D(\Omega)}^{\|\cdot\|_{k,p}}.\]

Quando $p=2$, escrevemos $H_0^k(\Omega)$ no lugar de $W^{k,2}_0(\Omega)$. 

\begin{exer}
Sejam \(1<{\color{blue} p }<+\infty\),  \(\Omega\) um aberto de \(\rd\) e \(u\in W_0^{1,{\color{blue} p }}(\Omega)\). Mostre que \(\tilde{u}\in W^{1,{\color{blue} p }}(\rd)\), onde  \(\tilde{u}\) é a extensão de \(u\) a zero fora de \(\Omega\). Além disso, 
\[
\partial_i \tilde{u}=\widetilde{\partial_i u},\ i=1,\ldots, \dm.
\]
\end{exer}
\subsection*{Exercício \theexer}
\end{frame}

\begin{frame}

\begin{prop}
 Sejam $1\leq p<+\infty$ e $k\geq 0$ um inteiro. Então $W_0^{k,p}(\R^\dm)=W^{k,p}(\R^\dm)$.
\end{prop}

\begin{lemaaux}
	Sejam $\rho\in C_c^\infty(\R^\dm)$ e $v\in W^{1,p}(\R^\dm)$, $1\leq p< +\infty$. Então 
	\[\rho\ast v\in C^\infty(\R^\dm) \cap W^{1,p}(\R^\dm), \text{ e } \partial_i(\rho\ast v)=\rho\ast \partial_i v,\  \forall i=1,\ldots,N.\]
\end{lemaaux}
\end{frame}
